\documentclass[11pt,a4paper]{article}
\usepackage[margin=2.5cm]{geometry}
\usepackage{amsmath}
\usepackage{booktabs}
\usepackage{array}
\usepackage{graphicx}
\usepackage{hyperref}
\usepackage{parskip}
\usepackage{xcolor}
\usepackage{caption}
\usepackage{subcaption}
\usepackage{float}

\graphicspath{{figures/}}

% Placeholder box for missing figures
\newcommand{\figplaceholder}[2]{%
  \begin{center}
    \fbox{\parbox[c][#1][c]{0.85\linewidth}{\centering\color{gray}
      \textbf{[FIGURE NEEDED]} \\ #2}}
  \end{center}
}

\title{\textbf{Gas Injury Detection} \\[0.4em]
       \large Wearable Sensor Classification --- Progress Report}
\author{}
\date{\today}

\begin{document}
\maketitle

% ─────────────────────────────────────────────────────────────────────────────
\section*{Key Results}

A Random Forest classifier was trained on four labeled sessions (two sweat,
two blood) and validated on a held-out 20\% test split.  The model was then
applied to a previously unseen sweat session to assess real-world
generalizability.

\begin{center}
\begin{tabular}{lrrr}
\toprule
Class    & Precision & Recall & Samples \\
\midrule
Baseline & 1.000     & 1.000  & 10 \\
Sweat    & 0.900     & 1.000  & 9  \\
Blood    & 1.000     & 0.900  & 10 \\
\midrule
\textbf{Overall accuracy} & \multicolumn{3}{c}{\textbf{0.967}} \\
\bottomrule
\end{tabular}
\captionof{table}{Run 2 validation results (80/20 stratified split, 29 test windows).
  \textbf{Precision} --- of all windows predicted as that class, how many were correct.
  \textbf{Recall} --- of all actual windows of that class, how many were detected.}
\end{center}

\begin{figure}[H]
  \centering
  \includegraphics[width=\linewidth]{sweat0_output}
  \caption{Sweat\_0 (Apr 04) model predictions. Color bands show predicted
    class per 60\,s window; opacity scales with confidence.
    Sweat was introduced at $\approx$3400\,s.}
\end{figure}

% ─────────────────────────────────────────────────────────────────────────────
\section{System Overview}

\begin{figure}[H]
  \centering
  \includegraphics[width=0.7\linewidth]{hardware}
  \caption{Sensor hardware --- ESP32 with MEMS gas sensor array and SPEC electrochemical sensors.}
\end{figure}

The system uses an ESP32 microcontroller logging readings at 1\,Hz over serial.
Two sensor families are integrated:

\begin{center}
\begin{tabular}{lll}
\toprule
Sensor type & Channels & Target analytes \\
\midrule
MEMS (metal-oxide) & VOC, NH3, HCHO & Volatile organic compounds \\
SPEC (electrochemical) & H2S, EtOH & Hydrogen sulfide, Ethanol \\
\midrule
Environmental & Temp, RH & Humidity compensation \\
\bottomrule
\end{tabular}
\end{center}

% ─────────────────────────────────────────────────────────────────────────────
\section{Data Collection}

\begin{center}
\begin{tabular}{llll}
\toprule
Session  & Date    & Sample type & Notes \\
\midrule
Sweat 1a & Apr 05  & Sweat (samples 2, 3) & Introduced $\approx$3850\,s \\
Sweat 1b & Apr 05  & Sweat (sample 1)     & After recording reset at 4000\,s \\
Blood 0  & Apr 06  & 1.5\,ml blood        & Introduced $\approx$1950\,s \\
Blood 1  & Apr 06  & 1.5\,ml blood        & Second run, introduced $\approx$2450\,s \\
\midrule
Sweat 0  & Apr 04  & Sweat (validation only) & NH3 pin miswired --- excluded from training \\
\bottomrule
\end{tabular}
\end{center}

\begin{figure}[H]
  \centering
  \begin{subfigure}[t]{0.49\linewidth}
    \includegraphics[width=\linewidth]{blood0_mems_raw}
    \caption{Blood 0 --- raw MEMS signals}
  \end{subfigure}
  \hfill
  \begin{subfigure}[t]{0.49\linewidth}
    \includegraphics[width=\linewidth]{blood0_mems_filtered}
    \caption{Blood 0 --- after filtering and drift correction}
  \end{subfigure}
  \caption{Example MEMS signal (VOC, NH3, HCHO) before and after the processing pipeline.}
\end{figure}

\begin{figure}[H]
  \centering
  \begin{subfigure}[t]{0.49\linewidth}
    \includegraphics[width=\linewidth]{blood0_spec_evn_raw}
    \caption{Blood 0 --- raw SPEC signals}
  \end{subfigure}
  \hfill
  \begin{subfigure}[t]{0.49\linewidth}
    \includegraphics[width=\linewidth]{blood0_spec_filtered}
    \caption{Blood 0 --- after filtering and stabilization offset}
  \end{subfigure}
  \caption{Example SPEC signal (H2S, EtOH) before and after the processing pipeline.}
\end{figure}

% ─────────────────────────────────────────────────────────────────────────────
\section{Processing Pipeline}

Two independent pipelines clean the raw data before modeling.

\subsection{MEMS Sensors}
\begin{enumerate}
  \item \textbf{Artifact removal} --- known electrical glitches patched by
        linear interpolation over flagged time windows.
  \item \textbf{Causal EMA smoothing} --- span 30\,s (80\,s for noisy NH3
        in Sweat 1b).  Causal to match real-time inference behavior.
  \item \textbf{Linear drift correction} --- a linear trend fitted over the
        pre-sample baseline window is subtracted from the full session,
        removing slow resistive sensor drift.
\end{enumerate}

\subsection{SPEC Sensors}
\begin{enumerate}
  \item \textbf{PPM conversion} from raw voltage using TIA gain and
        manufacturer sensitivity code.
  \item \textbf{Outlier removal} --- rolling IQR and rate-of-change detectors;
        outliers replaced by linear interpolation.
  \item \textbf{Three-stage smoothing} --- Savitzky--Golay pre-filter followed
        by causal Butterworth low-pass ($f_c = 0.005$\,Hz).
  \item \textbf{Boot-up stabilization offset} --- signal zeroed at the
        automatically detected stabilization point.
\end{enumerate}

% ─────────────────────────────────────────────────────────────────────────────
\section{Classification Model}

\subsection{Feature Engineering}

Non-overlapping 60-second windows are extracted from labelled regions.
Within each window, four aggregates are computed per channel:
mean, standard deviation, maximum, and linear slope.
Rate-of-change (ROC) and multi-scale rolling ROC columns (15\,s / 30\,s / 60\,s)
are included as features.  Humidity (\texttt{rh\_pct}) helper columns are
included; temperature is excluded.

\begin{center}
\begin{tabular}{lrr}
\toprule
& Run 1 & Run 2 \\
\midrule
Trees            & 200   & 100 \\
Feature columns  & 284   & 116 \\
Feature set      & All helpers & ROC only (no roll\_mean / roll\_std / acc) \\
Validation       & LOSO GroupKFold & 80/20 stratified split \\
\bottomrule
\end{tabular}
\end{center}

\subsection{Run 1 --- Leave-One-Session-Out (LOSO) GroupKFold}

Run 1 used \textbf{Leave-One-Session-Out (LOSO)} cross-validation.
With four sessions (Sweat 1a, Sweat 1b, Blood 0, Blood 1), the data is split
into four folds.  In each fold, one session is held back entirely for
validation and the model is trained on the remaining three.  This rotates
until every session has been the validation set exactly once.

\begin{center}
\begin{tabular}{llc}
\toprule
Fold & Train sessions & Validate on \\
\midrule
1 & Sweat 1b, Blood 0, Blood 1 & Sweat 1a \\
2 & Sweat 1a, Blood 0, Blood 1 & Sweat 1b \\
3 & Sweat 1a, Sweat 1b, Blood 1 & Blood 0 \\
4 & Sweat 1a, Sweat 1b, Blood 0 & Blood 1 \\
\bottomrule
\end{tabular}
\end{center}

Run 1 mean validation accuracy was \textbf{0.752} ($\sigma = 0.253$).
Fold 3 collapsed to 0.333 --- with only two blood sessions available,
holding one out leaves the model with a single blood reference session,
which is insufficient to learn a robust blood boundary.

\subsection{Run 2 --- 80/20 Stratified Split}

Run 2 used a simpler \textbf{80/20 random split} stratified by class label,
using all available data for training.  80\% of windows (across all sessions)
are used for training and 20\% for validation, with the split ensuring each
class is proportionally represented in both sets.  This makes full use of the
limited dataset and gives a reliable estimate of in-distribution performance.
True out-of-sample validation is deferred to future sessions (e.g.\ Sweat 0).

Run 2 overall accuracy: \textbf{0.967} across 29 test windows.

\subsection{Class Balancing}

Baseline windows outnumber labelled sample windows.  Baseline is randomly
undersampled to match the size of the larger sample class before training.

% ─────────────────────────────────────────────────────────────────────────────
\section{Validation on Unseen Session (Sweat 0)}

The trained model was applied to the April 04 sweat session, which was
never used during training.  The raw CSV was processed with the same
pipeline (resampled to 1\,Hz, drift corrected, offset corrected) and
classified window-by-window.

\begin{figure}[H]
  \centering
  \includegraphics[width=\linewidth]{sweat0_output}
  \caption{Sweat\_0 window-by-window predictions. Grey = baseline, blue = sweat, red = blood.
    Opacity reflects model confidence. Sweat introduced at $\approx$3400\,s (orange dashed line).}
\end{figure}

The model predicts baseline during the warmup phase and transitions to sweat near the
sample introduction point at 3400\,s, with confidence increasing as the signal stabilizes.
Early windows (0--480\,s) show mixed baseline/blood probability due to the initial
sensor transient before drift correction fully takes effect.

\begin{center}
\small
\begin{tabular}{rrrll}
\toprule
Window & $t_{\text{start}}$ (s) & $t_{\text{end}}$ (s) & Prediction & Confidence \\
\midrule
 1 &    0 &   59 & Baseline &  45\% \\
 2 &   60 &  119 & Baseline &  53\% \\
 3 &  120 &  179 & Baseline &  53\% \\
 4 &  180 &  239 & Baseline &  57\% \\
 5 &  240 &  299 & Baseline &  57\% \\
 6 &  300 &  359 & Baseline &  66\% \\
 7 &  360 &  419 & Baseline &  54\% \\
 8 &  420 &  479 & Baseline &  52\% \\
 9 &  480 &  539 & Sweat    &  42\% \\
10 &  540 &  599 & Sweat    &  41\% \\
11 &  600 &  659 & Baseline &  43\% \\
12 &  660 &  719 & Baseline &  46\% \\
13 &  720 &  779 & Baseline &  39\% \\
\multicolumn{5}{c}{$\vdots$} \\
65 & 3840 & 3899 & Sweat    &  44\% \\
66 & 3900 & 3959 & Sweat    &  46\% \\
\bottomrule
\end{tabular}
\captionof{table}{Sweat\_0 window predictions (selected rows). Full output: 66 windows over 3960\,s.}
\end{center}

% ─────────────────────────────────────────────────────────────────────────────
\section{Next Steps}

\begin{itemize}
  \item Collect additional labelled sessions (more subjects, varying
        sample volumes) to improve generalization.
  \item Collect more blood sessions to improve within-class coverage
        (currently only two blood sessions available).
  \item Implement the causal pipeline on the ESP32 for real-time
        classification output.
  \item Add humidity correction to the SPEC sensor pipeline --- electrochemical
        sensors are sensitive to relative humidity, and correcting for RH
        variation may improve signal consistency across sessions.
\end{itemize}

\end{document}
